%-----------------------------------------------------------------------------%
\chapter{\babSatu}
\label{bab:1}
%-----------------------------------------------------------------------------%
In this chapter, we will introduce the background that we will discuss in this thesis.

%-----------------------------------------------------------------------------%
\section{Background}
\label{sec:latarBelakang}
%-----------------------------------------------------------------------------%
A language model is a type of machine learning model that predicts sequences in natural language. A language model trained on a vast corpus of data with billions of parameters is called a large language model (LLM). Today, LLMs are widely utilized for tasks ranging from simple task such as translation to code writing and mathematics. 

In recent years, large language models have gained popularity due to their surprising effectiveness in solving complex reasoning problems. This capability was largely enabled by the emergence of "Chain-of-Thought" (CoT) prompting \citep{DBLP:conf/nips/Wei0SBIXCLZ22}, which explicitly encourages models to decompose multi-step problems into intermediate logical segments. By scaling parameters and training data, these models have developed abilities that allow them to simulate deductive paths, mimicking human-like deliberation rather than simply mapping inputs directly to outputs.

However, despite these advances in reasoning, large language models remain prone to hallucinations. Because they are fundamentally probabilistic autoregressive systems, LLMs predict the "most likely" next token based on surface-level statistical patterns rather than a grounded understanding of formal logic. Consequently, an error in an earlier reasoning steps can lead the rest of the reasoning steps astray, leading the model to confidently assert logically impossible or factually incorrect conclusions when tasked with complex proof planning.

To address these ungrounded hallucinations, researchers proposed Neurosymbolic AI \citep{Wan2024TowardsCA}, an approach that integrates the linguistic flexibility of LLMs with the absolute precision of symbolic logic solvers. One framework that popularizes this approach is LINC (Logical Inference via Neurosymbolic Computation) \citep{olausson-etal-2023-linc}. LINC reformulates reasoning tasks as a neurosymbolic process. Instead of relying on the LLM to perform the deductive reasoning itself, LINC uses the LLM solely as a semantic parser to translate natural language premises into strict First-Order Logic (FOL) expressions. These expressions are then offloaded to an external, deterministic theorem prover (such as Prover9) that calculates the final truth value. While highly effective at eliminating deductive errors, LINC is strictly deductive; it assumes that every necessary fact is explicitly stated in the problem. When human contexts leave out "obvious" commonsense information, the prover inevitably fails because the logic formulation is incomplete.

Recognizing this limitation, a newer approach named ARGOS (Abductive Reasoning with Generalization Over Symbolics) \citep{cotnareanu2026balanced} was developed. ARGOS tackles the problem of missing commonsense information by striking a balance between neural and symbolic components. Rather than relying entirely on a static translation, ARGOS dynamically searches for missing premises (abductive reasoning). It uses the logical backbone provided by a SAT solver \citep{VanHarmelen2008-VANHOK} to guide the LLM in generating new, relevant commonsense rules. These new rules are then scored and injected back into the logical solver in an iterative loop until the problem can be formally resolved, or it falls back to a neural baseline.

One of the most promising applications for this neurosymbolic approach is moral reasoning. Moral judgments are inherently subjective and depend heavily on cultural contexts and prevailing commonsense norms; thus, a single ethical dilemma may elicit entirely different conclusions across various demographics \citep{Bentahila2021}. Given this variability, it is critical that AI systems not only reach a moral decision but also transparently explain the underlying rationale. Therefore, evaluating a framework like ARGOS on complex moral dilemmas provides a rigorous testbed for its ability to navigate highly contextual, human-centric reasoning.

%-----------------------------------------------------------------------------%
\section{Problem Statement}
\label{sec:masalah}
%-----------------------------------------------------------------------------%
Although neurosymbolic AI has shown promise in deductive and relational reasoning tasks (e.g., CLUTRR \citep{sinha-etal-2019-clutrr}, ProntoQA \citep{PrOntoQA}), applying these rigid symbolic solvers to the nuanced and open-textured domain of moral reasoning introduces new challenges. Traditional datasets for neurosymbolic evaluation rely on strict, quantifiable First-Order Logic where missing commonsense rules can be systematically generated by linking explicit entities. In contrast, moral reasoning, specifically within the ExplainEthics dataset \citep{quan-etal-2024-enhancing}, requires evaluating highly subjective human norms that are difficult to logically quantify, as the scenarios lack explicit entities or variables (e.g., $x$ and $y$ in a relation $P(x,y)$). While these nuanced scenarios do not necessarily require us to use a specific logical formalism, we choose to represent them using Propositional Logic for the sake of simplicity. It remains unclear how a neurosymbolic framework like ARGOS must be adapted to function within moral reasoning tasks, and whether this adaptation provides a meaningful performance improvement over purely neural baseline methods when evaluating moral norm violations.

%-----------------------------------------------------------------------------%
\subsection{Research Questions}
\label{sec:definisiMasalah}
%-----------------------------------------------------------------------------%
In this thesis, we would like to answer these two questions:
\begin{enumerate}
	\item RQ1: How is ARGOS implemented and adapted to moral reasoning tasks?
	\item RQ2: How does the adapted ARGOS pipeline perform compared to baseline Chain-of-Thought (CoT) method on the ExplainEthics dataset?
\end{enumerate}

%-----------------------------------------------------------------------------%
\subsection{Scope and Limitations}
\label{sec:batasanMasalah}
%-----------------------------------------------------------------------------%
Due to time and resource constraints, we bound this research by the following assumptions and limitations:
\begin{enumerate}
	\item The experiment is conducted exclusively on the ExplainEthics dataset \citep{quan-etal-2024-enhancing}.
	\item The LLMs used for premise generation and the CoT fallback are limited to specific sizes across different model families (specifically Llama-3.1-8B-Instruct, Qwen/Qwen2.5-7B-Instruct, Qwen2.5-Coder-7B-Instruct, and Mistral-7B-Instruct) due to computational constraints.
\end{enumerate}

%-----------------------------------------------------------------------------%
\section{Research Purpose}
\label{sec:tujuan}
%-----------------------------------------------------------------------------%
The primary objectives of this research are:
\begin{enumerate}
	\item To adapt and implement the ARGOS neurosymbolic framework to operate within the highly subjective domain of moral reasoning.
	\item To empirically evaluate the performance and identify the fundamental bottlenecks of the adapted ARGOS pipeline compared to a baseline Chain-of-Thought (CoT) method on the ExplainEthics dataset.
\end{enumerate}

%-----------------------------------------------------------------------------%
\section{Thesis Organization}
\label{sec:sistematikaPenulisan}
%-----------------------------------------------------------------------------%
The organization of this thesis is as follows:
\begin{itemize}
	\item Chapter 1 \babSatu \\
	    This chapter covers the background, problem statement, research questions, scope, and purpose of the study.
	\item Chapter 2 \babDua \\
	    This chapter discusses the theoretical foundations of propositional logic, large language models, SAT solvers, and the ARGOS pipeline.
	\item Chapter 3 \babTiga \\
	    This chapter describes the research methodology, including the experimental design, dataset formulation, baseline models, and evaluation metrics.
	\item Chapter 4 \babEmpat \\
	    This chapter presents the experimental setup and the quantitative results obtained from running the pipeline.
	\item Chapter 5 \babLima \\
	    This chapter provides an analysis and discussion of the results with respect to the three research questions.
	\item Chapter 6 \kesimpulan \\
	    This chapter concludes the results of this thesis and offers suggestions for future research.
\end{itemize}
